\begin{textsample}{POS Dim 1 – 2015\_c5 – Score 10.00 – 2015\_c5/enade\_2015\_psico\_e312af574b.txt}  \label{ex:f1_pos_094}
CAPÍTULO 5 DISTRIBUIÇÃO DOS CONCEITOS 5.1 \textbf{PANORAMA} NACIONAL DA DISTRIBUIÇÃO DOS CONCEITOS O Conceito Enade é um indicador de qualidade que avalia o desempenho dos estudantes a partir dos resultados do Enade.
Os valores possíveis vão de 1 ( pior situação ) a 5 ( melhor situação ).
As Unidades de Observação com apenas um ou sem nenhum concluinte participante não obtêm o Conceito Enade, ficando Sem Conceito ( SC ).
A Tabela 5.1 apresenta a quantidade e distribuição de cursos de Psicologia participantes do Enade/2015, por faixa de conceito e Grande Região.
Mais uma vez, ressaltando que a diferença entre os cursos \textbf{tabulados} neste capítulo e no capítulo 2 corresponde aos cursos sem conceito, em princípio, aqueles sem alunos concluintes que participassem da prova ou que tivessem tão somente um aluno concluinte18.
Observando -se os dados da Tabela 5.1, nota -se que, dos 466 cursos participantes, 241 ( 51,7 \% ) classificaram -se com conceito 3, o valor modal.
Este foi também o conceito modal em mais quatro regiões : Nordeste ( 49,4 \% ), Sudeste ( 55,5 \% ), Sul ( 55,0 \% ) e Centro-Oeste ( 38,5 ).
Na região Norte, a moda foi o conceito 2 ( 48,4 \% ).
O conceito 4 foi o segundo mais frequente em nível nacional ( 21,2 \%, correspondendo a 99 cursos ) e o conceito 2, o terceiro ( 20,2 \%, correspondendo a 94 cursos ).
Houve, ainda, 23 cursos ( 4,9 \% ) que receberam conceito 5 e quatro ( 0,9 \% ) cursos que receberam conceito 1.
Dos 466 cursos de Psicologia, cinco cursos ( 1,1 \% ) ficaram sem conceito ( SC ).
O Gráfico 5.1 apresenta a distribuição do Conceito Enade segundo Grande Região : \textbf{poligonais} mais à esquerda representam Grandes Regiões com pior distribuição, e \textbf{poligonais} mais à direita, Grandes Regiões com melhores conceitos.
A região Norte ( linha azul ) apresenta os cursos em IES com conceitos mais baixos, ao passo que a região Nordeste ( linha verde ) apresenta as distribuições com valores maiores.
Cursos em IES das regiões Centro-Oeste ( linha preta ), Sul ( linha \textbf{roxa} ) e Sudeste ( linha vermelha ) apresentam uma situação intermediária. 18 Estes cursos com somente um concluinte não foram considerados no capítulo 2.
A região Norte participou com 31 cursos ou 6,7 \% do total nacional.
Nesta região, o conceito 2 foi atribuído a 15 cursos dos 31 participantes, o que \textbf{equivale} a 48,4 \% do total regional, além de ser o conceito modal na região.
Nenhum curso ficou sem conceito ( SC ) ou com conceito 5.
Um curso ( 3,2 \% ) foi avaliado com conceito 1, 12 ( 38,7 \% ) obtiveram conceito 3 e três cursos ( 9,7 \% ) receberam conceito 4.
A região Nordeste participou com 85 cursos ou 18,2 \% do total nacional.
Nesta região, 42 cursos, 49,4 \% em termos regionais, obtiveram conceito 3, o conceito modal para a região.
Os conceitos 2 e 4 foram atribuídos a, respectivamente, 19 e 14 cursos ( 22,4 \% e 16,5 \% ).
Nove cursos ( 10,6 \% ) obtiveram o conceito 5.
Apenas um curso ficou sem conceito ( SC ) e nenhum obteve conceito 1 nessa região.
Dos 200 cursos participantes da região Sudeste ( 42,9 \% do total nacional ), o conceito 3 também foi o conceito modal, tenho sido obtido por 111 cursos ( 55,5 \% ).
O conceito 1 foi atribuído a dois cursos ( 1,0 \% ) e os conceitos 2 e 4, a 32 e 48 cursos, correspondendo a 16,0 \% e 24,0 \%, respectivamente.
O conceito 5 foi atribuído a sete cursos ( 3,5 \% ).
Nenhum curso ficou sem conceito ( SC ).
Dos 111 cursos da região Sul, há predominância do conceito 3, conceito modal, conferido a 55,0 \%, ou seja, a 61 dos 111 cursos participantes na região Sul.
O conceito 2 foi atribuído a 18 cursos ( 16,2 \% ) e o conceito 4, a 23 cursos ( 20,7 \% ).
O conceito 5 foi recebido por cinco cursos ( 4,5 \% ), e nenhum curso ficou com conceito 1.
Na região Sul, quatro cursos ( 3,6 \% ) ficaram sem conceito ( SC ).
Já dos 39 cursos participantes na região Centro-Oeste, 15 receberam conceito 3, o conceito modal ( 38,5 \% ).
Nesta região, nenhum curso ficou sem conceito ( SC ).
Já um curso foi alocado no conceito 1, correspondendo a 2,6 \% dos cursos nessa região.
Dez cursos ( 25,6 \% ) foram avaliados com conceito 2, 11 cursos ( 28,2 \% ) obtiveram conceito 4 e somente dois cursos ficaram com conceito 5, correspondendo 5,1 \%, nessa região.
O Gráfico 5.1 apresenta a distribuição do Conceito Enade segundo Grande Região : \textbf{poligonais} mais à esquerda representam Grandes Regiões com pior distribuição, e \textbf{poligonais} mais à direita, Grandes Regiões com melhores conceitos.
A região Norte ( linha azul ) apresenta os cursos em IES com conceitos mais baixos, ao passo que a região Nordeste ( linha verde ) apresenta as distribuições com valores maiores.
Cursos em IES das regiões Centro-Oeste ( linha preta ), Sul ( linha \textbf{roxa} ) e Sudeste ( linha vermelha ) apresentam uma situação intermediária. 5.2 CONCEITOS POR CATEGORIA ADMINISTRATIVA E POR GRANDE REGIÃO A Tabela 5.2 apresenta a distribuição dos cursos participantes do Enade/2015 de Psicologia, por Categoria Administrativa, de acordo com os conceitos por eles alcançados, segundo as Grandes Regiões brasileiras.
Dos 466 cursos participantes, 85 ( 18,2 \% ) eram \textbf{ministrados} em Instituições \textbf{Públicas} e 381 ( 81,8 \% ), em Instituições Privadas.
De acordo com as informações da Tabela 5.2, em termos nacionais, entre Instituições \textbf{Públicas} e Privadas, nota -se que das 23 IES que receberam o conceito 5, 21 eram Públicas e duas Privadas.
Dos 85 cursos participantes de IES Públicas, o conceito 4 foi o valor modal, atribuído a 35 cursos ( 41,2 \% ).
Entre os demais cursos de Instituições \textbf{Públicas} participantes, dois cursos ( 2,4 \% ) foram avaliados com conceito 1, oito cursos ( 9,4 \% da categoria ) obtiveram conceito 2, 18 cursos ( 21,2 \% ) receberam conceito 3 e, como já comentado, 21 ( 24,7 \% ) foram alocados ao conceito 5.
Nesta categoria, apenas um curso ficou sem conceito ( SC ).
Na rede privada, o conceito modal foi 3, atribuído a 223 cursos dos 381 da categoria.
Entre os demais cursos participantes, dois receberam conceito 1 e 86 o conceito 2.
O conceito 4 foi atribuído a 64 cursos, e o conceito 5, a dois cursos.
Nesta Categoria Administrativa, outros quatro cursos ficaram sem conceito ( SC ).
Na análise por região, observa -se que, na região Norte, as Instituições Privadas participaram com 25 cursos ( 80,6 \% do total regional ), dos quais nenhum ficou sem conceito ( SC ) ou com conceitos 4 ou 5.
O conceito modal para as Instituições Privadas na região foi 2, com 13 cursos, correspondendo a 52,0 \%.
Um curso recebeu conceito 1 Na região Norte, as Instituições Privadas participaram com 25 cursos ( 80,6 \% do total regional ), dos quais nenhum ficou sem conceito ( SC ) ou com conceitos 4 ou 5.
O conceito modal para as Instituições Privadas na região foi 2, com 13 cursos, correspondendo a 52,0 \%.
Um curso recebeu conceito 1 e 11 cursos foram alocados no conceito 3.
As Instituições Públicas participaram com seis cursos na região Norte ( 19,4 \% do total regional ), dos quais três ( 50,0 \% ) obtiveram conceito 4, o valor modal, dois cursos ( 33,3 \% ) receberam conceito 2 e um único curso ( 16,7 \% ) obteve conceito 3.
Nenhum curso ficou sem conceito ( SC ) ou recebeu os conceitos 1 ou 5 nesta combinação de categoria e região.
Na região Nordeste, a rede privada concentrou 64 dos 85 cursos participantes, o equivalente a 75,3 \% do total da região.
Dentre os cursos oferecidos por IES Privadas no Nordeste, 40 cursos ( 62,5 \% ) ficaram alocados no conceito 3, conceito modal, 18 cursos ( 28,1 \% ) obtiveram conceito 2, quatro ( 6,3 \% ) receberam conceito 4 e um curso ( 1,6 \% ) ficou com conceito 5.
Apenas um curso ficou sem conceito ( SC ) enquanto nenhum foi avaliado com conceito 1.
As Instituições \textbf{Públicas} dessa região participaram com 21 cursos ( 24,7 \% ), dos quais dez ( 47,6 \% ) obtiveram conceito 4, o valor modal, e oito cursos ( 38,1 \% ) receberam conceito 5.
Um curso ( 4,8 \% ) foi alocado no conceito 2, enquanto dois ( 9,5 \% ) obtiveram conceito 3.
Nenhum curso ficou sem conceito ( SC ) ou obteve conceito 1 nesta combinação de categoria e região.
Na região Sudeste, a proporção de cursos da rede privada, 85,5 \%, foi mais elevada do que nas demais regiões brasileiras, correspondendo a 171 dos 200 cursos participantes.
Nesta categoria, na região Sudeste, o conceito modal foi 3 ( 105 cursos ).
Nesta combinação de Categoria Administrativa e Grande Região, nenhum curso ficou sem conceito ( SC ) ou recebeu conceito 1 ou 5.
Os demais foram avaliados com conceito 2 ( 29 cursos ) e conceito 4 ( 37 cursos ).
Entre os 29 cursos oferecidos por Instituições \textbf{Públicas} na região Sudeste, o conceito modal foi 4 ( 11 cursos ), e os demais cursos receberam os conceitos 1 ( dois cursos ), 2 ( três cursos ), 3 ( seis cursos ) e 5 ( sete cursos ).
Na região Sudeste, nesta categoria, nenhum curso ficou sem conceito ( SC ).
As Instituições Privadas concentraram 94 dos 111 cursos participantes da região Sul, 84,7 \% do total regional.
Desses, 55 ficaram com conceito 3, o conceito modal.
Nesta combinação de Categoria Administrativa e Grande Região, 18 cursos ficaram com conceito 4, 17 cursos tiraram conceito 2 e apenas um curso obteve conceito 5.
Nenhum curso ficou com conceito 1 e três cursos ficaram sem conceito ( SC ).
As Instituições \textbf{Públicas} na região Sul participaram com 17 cursos ( 15,3 \% ) e o conceito modal foi 3, com seis cursos.
Um curso obteve conceito 2, cinco cursos foram avaliados com conceito 4 e quatro cursos receberam conceito 5.
Apenas um curso ficou sem conceito ( SC ) enquanto nenhum recebeu conceito 1.
Na região Centro-Oeste, 27 dos 39 cursos participantes eram de Instituições Privadas ( 69,2 \% em termos regionais ).
Destes, 12 concentraram -se no conceito 3, conceito modal.
Entre os demais, um curso recebeu conceito 1, nove obtiveram conceito 2 e cinco receberam conceito 4.
Nenhum curso foi avaliado com conceito 5 ou ficou sem conceito ( SC ) nesta categoria.
Dos 12 cursos oferecidos por Instituições \textbf{Públicas}, seis foram avaliados com conceito 4, conceito modal, um curso obteve conceito 2, três cursos receberam conceito 3 e dois cursos foram avaliados com conceito 5.
Nesta região, nenhum curso de IES \textbf{pública} ficou sem conceito ( SC ) ou recebeu conceito 1.
O Gráfico 5.2 apresenta a distribuição do Conceito Enade segundo a Categoria Administrativa da IES.
Os cursos em IES Públicas ( linha azul ) apresentam uma \textbf{poligonal} mais à direita do que os em IES Privadas ( linha verde ) e, consequentemente, uma distribuição de Conceitos Enade com valores maiores. 5.3 CONCEITOS POR ORGANIZAÇÃO ACADÊMICA E POR GRANDE REGIÃO Na Tabela 5.3, encontra -se a distribuição dos conceitos atribuídos aos cursos participantes do Enade/2015 na Área de Psicologia, por Organização Acadêmica, segundo as Grandes Regiões brasileiras.
Dos 466 cursos de Psicologia participantes, 218 eram oferecidos em Universidades ; 67 cursos em Centros Universitários ; e 181 cursos em Faculdades.
Esta distribuição corresponde, respectivamente, a 46,8 \%, 14,4 \% e 38,8 \% dos cursos.
De acordo com os dados apresentados, de todos os 23 cursos que tiveram 5 como conceito, 21 eram vinculados a Universidades.
Quatro cursos oferecidos em Universidades ficaram sem conceito ( SC ).
Esse tipo de Organização Acadêmica teve o conceito 3 como modal, atribuído a 102 cursos ( 46,8 \% ).
Os demais cursos de Universidades avaliados receberam os conceitos 1 ( dois cursos ), 2 ( 19 cursos ), 4 ( 70 cursos ) e conceito 5 ( 21 cursos, como já mencionado ).
Entre os cursos em Centros Universitários, o conceito modal foi 3, atribuído a 41 cursos ( 61,2 \% ).
Neste tipo de Organização Acadêmica, 12 cursos obtiveram conceito 2 e 14 ficaram com conceito 4.
Nenhum curso ficou sem conceito ( SC ) ou foi avaliado com os conceitos 1 ou 5.
Nas Faculdades, 98 cursos ficaram alocados no conceito 3, conceito modal.
Dos demais, dois cursos obtiveram conceito 1, 63 cursos foram avaliados com o conceito 2, 15 cursos ficaram com o conceito 4 e dois cursos receberam conceito 5.
Apenas um dos 181 cursos ficou sem conceito ( SC ) neste tipo de Organização Acadêmica.
Considerando -se separadamente as regiões brasileiras, verifica -se que, na região Norte, dos 31 cursos participantes, seis eram oferecidos por Universidades, todos os quais avaliados com os conceitos 3 ou 4 ( três cursos cada ).
Sendo assim, nenhum curso ficou sem conceito ( SC ) ou recebeu os demais conceitos nesta categoria e região.
Os Centros Universitários da região Norte foram representados por quatro cursos.
Três cursos receberam conceito 3, conceito modal, e um curso foi avaliado com conceito 2.
Nenhum curso ficou sem conceito ( SC ) ou recebeu os demais conceitos.
As Faculdades participaram com 21 cursos na região Norte, que receberam conceitos 1 ( um curso ), 2 ( 14 cursos, conceito modal ) e 3 ( seis cursos ).
Nenhum curso ficou sem conceito ( SC ) ou recebeu os demais conceitos.
Na região Nordeste, as Universidades participaram com 26 dos 85 cursos na Área de Psicologia.
O conceito modal foi 4, atribuído a dez cursos.
Os conceitos 3 e 5 foram atribuídos a oito cursos cada.
Nenhum curso oferecido por Universidade no Nordeste ficou sem conceito ( SC ) ou recebeu os demais conceitos.
Os Centros Universitários \textbf{contaram} com dez cursos participantes na região Nordeste, sendo que cinco obtiveram conceito 3, conceito modal, três receberam conceito 2 e dois cursos tiveram conceito 4.
Nenhum curso ficou sem conceito ( SC ) ou obteve os demais conceitos.
As Faculdades foram representadas por 49 cursos na região Nordeste, dos quais 29 receberam conceito 3, conceito modal, 16 obtiveram conceito 2, dois cursos obtiveram conceito 4 e um curso foi avaliado com conceito 5.
Apenas um curso ficou sem conceito ( SC ), e nenhum recebeu conceito 1.
Na região Sudeste, as Universidades concentraram 106 dos 200 cursos da região.
Entre os cursos em Universidades, nessa região, o conceito modal foi 3, atribuído a 53 cursos, enquanto dois cursos ficaram com conceito 1, 12 obtiveram conceito 2, 32 receberam conceito 4 e sete cursos foram avaliados com o conceito 5.
Nenhum curso ficou sem conceito ( SC ).
Os Centros Universitários participaram com 35 cursos na região Sudeste, dos quais 23 obtiveram o conceito modal, 3, enquanto os conceitos 2 e 4 foram atribuídos a seis cursos cada.
Nenhum curso ficou sem conceito ( SC ) ou recebeu conceitos 1 ou 5.
As Faculdades foram representadas por 59 cursos na região Sudeste, dos quais 35 receberam o conceito 3, conceito modal.
Dos demais, 14 cursos receberam o conceito 2 e dez cursos obtiveram o conceito 4.
Nenhum curso ficou sem conceito ( SC ) ou recebeu os conceitos 1 ou 5 nessa categoria e região.
Dos 111 cursos da região Sul, 61 eram de Universidades, para os quais o conceito modal foi 3, atribuído a 30 cursos.
Nesse tipo de organização, quatro cursos ficaram sem conceito ( SC ).
Os demais receberam os conceitos 2 ( seis cursos ), 4 ( 17 cursos ), 5 ( quatro cursos ).
Nenhum curso recebeu conceito 1.
Dos 13 cursos participantes de Centros Universitários da região Sul, oito obtiveram conceito 3, conceito modal, dois receberam conceito 2 e três cursos ficaram com conceito 4.
Nenhum curso ficou sem conceito ( SC ) ou recebeu os conceitos 1 ou 5.
Foram 37 os cursos vinculados a Faculdades na região Sul, sendo que 23 receberam conceito 3, o conceito modal.
Dez cursos obtiveram conceito 2, três cursos receberam conceito 4 e um curso ficou com o conceito 5.
Nenhum curso ficou sem conceito ( SC ) ou recebeu o conceito 1 nessa categoria e região.
Na região Centro-Oeste, 19 dos 39 cursos eram de Universidades.
Nesse tipo de organização, nenhum curso ficou sem conceito ( SC ) ou com conceito 1.
Os conceitos foram distribuídos de forma \textbf{bimodal}, uma vez que os conceitos 3 e 4, que registraram as maiores frequências, foram obtidos por oito cursos cada.
Apenas um curso obteve conceito 2 e dois cursos foram avaliados com o conceito 5.
Os Centros Universitários da região Centro-Oeste \textbf{contaram} com cinco cursos, que obtiveram conceitos 3 ( dois cursos ) e 4, conceito modal ( três cursos ).
Nenhum curso ficou sem conceito ( SC ) ou recebeu os conceitos \textbf{restantes} na categoria.
Dos 15 cursos oferecidos em Faculdades na região Centro-Oeste, nove ficaram alocados no conceito 2, conceito modal, um curso obteve conceito 1 e cinco cursos receberam conceito 3.
Nenhum curso ficou sem conceito ou recebeu os conceitos 4 ou 5.
O Gráfico 5.3 apresenta a distribuição do Conceito Enade segundo a Organização Acadêmica da IES.
Os cursos em Universidades ( linha azul ) apresentam uma \textbf{poligonal} mais à direita do que as demais, \textbf{denotando} uma melhor distribuição de Conceitos Enade.
No outro extremo, temos os cursos oferecidos em Faculdades ( linha vermelha ) com a \textbf{poligonal} mais à esquerda e, consequentemente, uma distribuição de Conceitos Enade com menores valores.
Cursos de Centros Universitários ( linha verde ) apresentam uma situação intermediária.

% matched lemmas: bimodal, contar, denotar, equivaler, ministrar, panorama, poligonal, pública, restante, roxa, tabular
\end{textsample}
